\documentclass[12pt,a4paper]{article}

\usepackage{cmap}
\usepackage[T2A]{fontenc}
\usepackage[utf8]{inputenc}
\usepackage[russian]{babel}
\usepackage{amsmath,amssymb}
\usepackage[a4paper,margin=2cm]{geometry}

\setlength{\parindent}{0pt}
\setlength{\parskip}{0.45em}
\sloppy

\begin{document}

\begin{center}
{\Large\bfseries Билет 9: вопросы для проверки знаний}
\end{center}

\noindent\fbox{%
  \begin{minipage}{\dimexpr\linewidth-2\fboxsep-2\fboxrule\relax}
  \normalfont
Несобственный интеграл. Критерий Коши сходимости интеграла. Интегралы от
знакопостоянных функций, признак сравнения. Интегралы от знакопеременных
функций, сходимость и абсолютная сходимость. Признаки Дирихле и Абеля
сходимости интегралов.
  \end{minipage}%
}

\section*{Вопросы на понимание}

\begin{enumerate}
\item Чем отличается несобственный интеграл по бесконечному промежутку от несобственного интеграла с особенностью в конечной точке?
\item Почему при нескольких особых точках интеграл разбивают на части и требуют сходимости каждой части отдельно?
\item Может ли изменение функции на отрезке, удаленном от особой точки, изменить сходимость несобственного интеграла? Почему?
\item Что в критерии Коши означает малость интегралов по хвостам между произвольными \(b_1\) и \(b_2\)?
\item Почему для неотрицательной функции расходимость связана с неограниченным ростом частичных интегралов, а не с колебаниями?
\item В первом признаке сравнения почему сходимость большей неотрицательной функции влечет сходимость меньшей, но не наоборот?
\item Почему во втором признаке сравнения достаточно сравнивать функции только около особой точки?
\item Чем условная сходимость отличается от абсолютной на уровне компенсации положительных и отрицательных значений функции?
\item Может ли интеграл от знакопеременной функции сходиться, если интеграл от ее модуля расходится? Какой тип сходимости это иллюстрирует?
\item В признаке Дирихле какую роль играет ограниченность первообразной функции \(f\)?
\item Почему в признаке Дирихле условие \(g(x)\to0\) нельзя заменить одной монотонностью функции \(g\)?
\item Как в признаке Абеля вычитание предела функции \(g\) сводит задачу к идее признака Дирихле?
\end{enumerate}

\section*{Источники для сверки}

\texttt{target/answer\_question\_09.tex};
\texttt{IvGE\_1.pdf}, стр. 244--250, 252--262.

\end{document}
